\documentclass[11pt]{article}
\usepackage[T1]{fontenc}
\usepackage[utf8]{inputenc}
\usepackage{amsmath,amssymb}
\usepackage[margin=2.5cm]{geometry}
\usepackage{listings}
\lstset{basicstyle=\ttfamily\small,columns=fullflexible,frame=single,
        framerule=0.3pt,xleftmargin=1em,xrightmargin=1em,showstringspaces=false,
        literate={ε}{{$\varepsilon$}}1}
\newcommand{\m}{\texttt{|}}
\title{A Tutorial on Unary Arithmetic\\\large Numbers as strings of marks, no integers required}
\author{Companion to \emph{From Unary to Binary}}
\date{\today}
\begin{document}
\maketitle

\section{Motivation}
The companion paper starts from a featureless infinity. Splitting it cannot give
the parts any intrinsic property, so each part is again the same kind of infinity,
and the parts are interchangeable: unary objects. A collection of $n$ of them is
characterised by $n$ alone and is invariant under $\mathrm{Sym}(n)$, so it has no
order and no arithmetic of its own. Nested splits address the parts by their
position in the nesting, and reading such an address as a string over $\{0,1\}$ is
binary information (Section~\ref{sec:bin}). All of this, the arithmetic included,
is a reading made by rich structures such as ourselves, not a feature of the base.
This tutorial shows what arithmetic such a reader can carry out on unary numerals
by string operations alone, and how binary relates to it.

\section{The idea}
A \emph{unary numeral} is a string of identical marks. We write the mark as \m{}, so
\m\m\m{} is the number three and the empty string $\varepsilon$ is zero. There are no
place values and no other symbols. The only thing that distinguishes one numeral
from another is how many marks it has.

Read as a structure, a numeral is a chain of nested parts, each extending the
previous one by one mark. This is more than a bare collection, but since all marks
are identical, only the length of the chain matters. Two remarks on scope:
\begin{itemize}
  \item The operations are conventions applied by a reader. A bare collection of
        unary digits has no arithmetic of its own. The code in \texttt{unary.py} is
        one such reader.
  \item The chain gives only the successor (append a mark). Addition,
        multiplication and the rest are defined on top of it, as string operations.
\end{itemize}
The rule of the game is that no operation may use integer arithmetic. The single
exception is \texttt{to\_int}, a reading taken from outside so that we can display
and test the answer. No operation uses it, and a test enforces this.

\section{Reading, writing, and order}
\begin{lstlisting}
>>> F = Unary.from_int
>>> F(3), F(0)
(|||, ε)
\end{lstlisting}
\textbf{Order is a prefix test.} Because every numeral is a run of marks,
$a\le b$ exactly when $a$ is a prefix of $b$. Only matching from the left is needed.
\begin{lstlisting}
def __le__(self, o): return o.s.startswith(self.s)
\end{lstlisting}
\textbf{Successor} appends one mark. Everything else is built from this.

\section{Addition: concatenation}
\[
  \m\m\m \;+\; \m\m\m\m \;=\; \m\m\m\m\m\m\m
\]
\begin{lstlisting}
def __add__(self, o): return Unary(self.s + o.s)
\end{lstlisting}
Associativity and the identity $\varepsilon$ are properties of concatenation.

\paragraph{Why addition commutes here.} For strings over a two-letter alphabet,
concatenation does \emph{not} commute: $01\cdot 10 \neq 10\cdot 01$. With a one-letter
alphabet every string is a run of the same mark, so $a\cdot b = b\cdot a$
automatically. Commutativity of unary addition is therefore a consequence of having
no distinctions between marks, which echoes the interchangeability of the unary
objects in the paper.

\section{Subtraction: removing a prefix}
$a-b$ is defined when $b\le a$ and removes $b$ from $a$. (Since both are runs of
identical marks, removing from the front or the end gives the same result.)
\begin{lstlisting}
def __sub__(self, o):
    if not o <= self: raise ValueError("a - b undefined for b > a ...")
    return Unary(self.s.replace(o.s, "", 1))
\end{lstlisting}
\begin{lstlisting}
>>> F(5) - F(2)
|||
>>> F(2) - F(3)
ValueError: a - b undefined for b > a in the naturals
\end{lstlisting}
Subtraction is partial because the naturals have no negatives. Nothing else is
needed.

\section{Multiplication: substitution}
Replace each mark of $a$ by a copy of $b$.
\[
  \m\m\m \times \m\m\m\m
  \;\Rightarrow\; (\m\m\m\m)(\m\m\m\m)(\m\m\m\m)
  \;=\; \m\m\m\m\m\m\m\m\m\m\m\m
\]
\begin{lstlisting}
def __mul__(self, o): return Unary(self.s.replace(D, o.s))
\end{lstlisting}
Multiplying by $\varepsilon$ erases every mark (the result is $\varepsilon$), and
multiplying by a single mark leaves the numeral unchanged. Commutativity is not
obvious from the rule, since ``$a$ copies of $b$'' and ``$b$ copies of $a$'' are
different recipes. The rectangle picture shows they give the same number of marks:
arrange $a$ rows of $b$ marks and read the columns instead of the rows.
Distributivity is the same picture cut into two blocks.

\section{Division, powers and the gcd}
\textbf{Division} is repeated subtraction: remove $b$ from $a$ until you cannot,
recording the removals with a growing numeral.
\begin{lstlisting}
>>> divmod(F(13), F(4))
(|||, |)
\end{lstlisting}
\textbf{Powers} are repeated multiplication starting from a single mark, so
$a^{\varepsilon}=\m$. The code gives $0^0=1$, the usual convention.
\begin{lstlisting}
>>> F(2) ** F(3)
||||||||
\end{lstlisting}
\textbf{Greatest common divisor} is Euclid's algorithm using only the remainder:
\begin{lstlisting}
>>> F(12).gcd(F(18))
||||||
\end{lstlisting}

\section{Converting to binary}\label{sec:bin}
Binary is a compression of unary. The conversion shows why.

\paragraph{Unary to binary: pair off the marks.}
Group the marks into pairs. A leftover single mark means the next bit is $1$,
no leftover means $0$. The pairs become a new, shorter numeral, and the process
repeats. Bits come out least significant first.
\[
\m\m\m\m\m\m\;\to\; \underbrace{\m\m}\,\underbrace{\m\m}\,\underbrace{\m\m}
\]
Six marks give three pairs and no leftover, so the bit is $0$. Three marks give one
pair and one leftover, so the bit is $1$. One mark gives no pair and one leftover,
so the bit is $1$ and nothing remains. Reading the bits from last to first gives
$110$, the binary form of six.
\begin{lstlisting}
>>> F(6).to_binary()
'110'
\end{lstlisting}
Pairing off is the inverse of the two-way split. It recovers two-ness from marks and
does not derive it.

\paragraph{Binary to unary: Horner's rule.} Start from $\varepsilon$. For each bit,
double the string (concatenate it with itself) and append a mark if the bit is $1$.
\begin{lstlisting}
>>> Unary.from_binary("1101")
|||||||||||||
\end{lstlisting}
This is thirteen marks. The doubling step is multiplication by two, so this is
Horner's method for evaluating $\sum b_i 2^i$.

\paragraph{Addresses and weights.} In the paper, nested splits give each part an
address: a path from the base, written as a string over $\{0,1\}$ whose symbols name
only which branch was taken. Nothing in the splits assigns a weight to a position.
\texttt{from\_binary} is a reader's convention that does: it maps an address to a
numeral, giving position $j$ the weight $2^j$, and \texttt{to\_binary} maps the
numeral back. The weights are on the reader's side, not in the digits or the splits.

\paragraph{The size difference.} A unary numeral for $n$ has $n$ marks. Its binary
form has about $\log_2 n$ bits. The conversion shrinks the string at the cost of a
second symbol and a positional convention, which a collection of unary digits does
not supply.

\section{What this does and does not show}
\begin{itemize}
  \item \textbf{Verified.} (\texttt{unary.py}, \texttt{test\_unary.py}) For all
        pairs $0\le a,b\le 24$ the string results for $+,-,\times$, $\div$, order and
        gcd agree with ordinary integer arithmetic (powers for small inputs), and
        binary conversion round-trips for $0\le n<200$. The usual algebraic laws
        (commutativity, associativity, distributivity, division with remainder) are
        checked directly, and a test confirms that no operation calls \texttt{to\_int}.
  \item \textbf{The machinery is not unary.} The host language implements
        \texttt{replace} and \texttt{startswith} using its own counting and a rich
        alphabet. The tutorial claims only that the \emph{numerals} and the
        \emph{rules on them} need no integers.
  \item \textbf{The converter uses a second symbol.} \texttt{to\_binary} marks pairs
        with \texttt{\#}, and its output is a string over $\{0,1\}$. This is the
        reader introducing a binary distinction on its side.
  \item \textbf{Scope.} This is primitive-recursive arithmetic with host-language
        loops. It does not show Turing-completeness, which the paper takes from the
        standard route.
  \item \textbf{Cost.} Unary arithmetic is slow and long: multiplication takes time
        proportional to $ab$, and $a^b$ has $a^b$ marks. That is the price of having
        no positional structure.
\end{itemize}

\section*{Exercises}
\begin{enumerate}
  \item Prove that $a\le b$ iff there is a $c$ with $a+c=b$, using only prefixes.
  \item Show $(a+b)\times c = a\times c + b\times c$ using substitution.
  \item Replace \texttt{to\_binary} by a version that never builds the symbol
        \texttt{\#}. Is it possible with only the original mark?
  \item Write \texttt{to\_ternary} using groups of three marks.
  \item Give a numeral-level argument that $\gcd(a,b)=\gcd(b,a \bmod b)$.
  \item Show that the $2^d$ addresses of length $d$ map under \texttt{from\_binary}
        to distinct numerals, and that \texttt{to\_binary} recovers each address up
        to leading zeros. Why are leading zeros lost?
\end{enumerate}
\end{document}
